\section{Introducción}

Este informe documenta el Trabajo práctico 2 del curso de Principios de
Inteligencia Artificial. El trabajo consiste en detectar mensajes spam del
dataset SMSSpamCollection \cite{almeida2011sms,uci_sms_spam}, un conjunto
público de mensajes SMS etiquetados como \emph{ham} o \emph{spam}, con tres
enfoques implementados en PyTorch \cite{paszke2019pytorch}:

\begin{enumerate}
    \item clasificación supervisada con regresión logística sobre las
    características de un codificador, con y sin reducción de
    dimensionalidad por análisis de componentes principales;
    \item detección de anomalías con densidades por histograma y un umbral
    de verosimilitud;
    \item clasificación con modelos grandes de lenguaje.
\end{enumerate}

La implementación completa está en el notebook
\texttt{tp2\_detector\_spam.ipynb}, organizado en celdas de definiciones,
configuración, experimentos y análisis. Las pruebas unitarias, entregadas
en la carpeta \texttt{tests/} junto al notebook, cargan las mismas
definiciones sin ejecutar descargas ni experimentos. Las métricas, figuras y
tablas de este informe provienen de la ejecución de las secciones del
notebook. El punto 3 evalúa una línea base sin ajuste (3.1), dos modelos
adicionales con cuantización NF4 (3.2) y el ajuste de Qwen 14B con LoRA en
una partición fija (3.3), que eleva el F1 de 0.7120 a 0.9708 en su prueba
reservada.

\subsection{Protocolo común}

\begin{itemize}
    \item \textbf{Etiquetas.} \emph{ham} $=0$ y \emph{spam} $=1$; spam es la
    clase positiva del F1-score y la anomalía del detector.
    \item \textbf{Métricas.} Tasa de aciertos y F1-score.
    \item \textbf{Particiones.} 30 particiones estratificadas de 70~\% de
    entrenamiento y 30~\% de prueba, una por semilla.
    \item \textbf{Muestreo de los modelos de lenguaje.} 30 muestras
    estratificadas del 30~\% del dataset.
    \item \textbf{Codificador.} \texttt{sentence-transformers/all-mpnet-base-v2},
    con vectores de 768 dimensiones (Sección~\ref{sec:caracteristicas}).
    \item \textbf{Hiperparámetros de la regresión logística.} Calibrados por
    F1 en una validación interna de cada corrida, tomada solo de su entrenamiento
    (Sección~\ref{sec:calibracion}).
    \item \textbf{Dimensión óptima del PCA.} Se eligió $D=384$ para la
    regresión logística, la mitad de las 768 características
    (Sección~\ref{sec:pca}).
    \item \textbf{Detector de anomalías.} Entre los casos de 1.5, el detector
    de histogramas usa $D=2$, elegida en una validación interna de los 30
    entrenamientos; las densidades se construyen solo con ham y el umbral es
    el cuantil 5~\% de un 20~\% de ham reservado (Sección~\ref{sec:punto2}).
    Esta variante difiere de reutilizar $D=384$, como pide la lectura literal
    del apartado 2.1; por eso también se informa el detector con $D=384$.
    \item Punto 3.1: \texttt{Qwen/Qwen2.5-3B-Instruct}, prompt sin ejemplos
    en inglés y decodificación voraz; las respuestas ambiguas cuentan como
    ham. Las 30 muestras estratificadas del 30~\% usan las semillas de los
    puntos 1 y 2 (Sección~\ref{sec:punto3}). Para 3.2 se usan
    Mistral-7B-Instruct-v0.3 y Qwen2.5-14B-Instruct, ambos con NF4 de 4 bits
    y cálculo bfloat16. En 3.3 se ajusta Qwen 14B con LoRA sobre NF4, con un
    entrenamiento y una prueba reservada (Sección~\ref{sec:punto33}).
\end{itemize}
